\documentclass{report}
\usepackage{amsmath,amssymb}
\setlength{\parindent}{0mm}
\setlength{\parskip}{1em}
\begin{document}
\begin{center}
\rule{6in}{1pt} \
{\large Laguna Ignacio \\
{\bf Protecting Scientific Applications From Silent Data Corruption via Static Analysis and Machine Learning}}

Lawrence Livermore National Laboratory \\ P O Box 808 \\ L-561 \\ Livermore \\ CA 94551-0808
\\
{\tt ilaguna@llnl.gov}\\
Martin Schulz\\
David F. Richards\\
Jon Calhoun\\
Luke Olson\end{center}

A possible increase of soft error rates is one of the major concerns for
the HPC community as larger and more complex HPC systems are built. In
exascale systems parallel scientific applications may be subjected to
higher degrees of silent data corruption (SDC) errors, which may arise
from unprotected components of the machine. As a result, programmers of
scientific codes would have to incorporate protection mechanisms in
computation code that could be vulnerable to these errors.
In this talk we present a framework to automatically protect scientific
applications from SDC in their output. The framework uses fault injection
and machine learning to learn code instructions that may yield SDC in
application�s output. Using the LLVM compiler, these vulnerable
instructions are duplicated to detect errors. We present case studies of
the technique in HPC kernels and mini applications, such as AMG and
HPCCG, and show that the overhead of the technique is small in comparison
to existing instruction duplication approaches.


\end{document}
